\section{Insolvability of the Quintic}

\exercise{56.1}{A radical extension in characteristic two}
The polynomial $x^2+x+1$ is irreducible over $\Z_2$ and its splitting field is $K=\operatorname{GF}(4)$. A square root of an element of $\Z_2$ cannot generate it: the only elements are 0,1, whose square roots in any extension are uniquely 0,1 in characteristic 2. Nevertheless $K$ is a radical extension. A root $\alpha$ of the polynomial satisfies $\alpha^3=1$, and $K=\Z_2(\alpha)$ is obtained by adjoining this cube root of the base-field element 1. A radical need not be a square root or the real positive choice of a root.

\exercise{56.2}{An even polynomial of degree eight}
Yes in characteristic zero. Substitute $y=x^2$ to get the quartic $ay^4+by^3+cy^2+dy+e$. Every quartic is solvable by radicals, so all four $y$-roots lie in a radical extension. Adjoining square roots of each of them, one after another, gives a further radical extension containing both square roots of every $y$-root. These are exactly the roots of the original polynomial. Thus its splitting field is contained in a radical extension.

\exercise{56.3}{Ten statements about radicals and solvability}
\begin{enumerate}[label=(\alph*)]
\item True. This is the definition of solvability by radicals.
\item True in characteristic zero, by the solvable-Galois-group criterion.
\item True. The roots of $x^{17}-5$ are products of a seventeenth root of 5 and powers of a seventeenth root of unity; adjoining these radicals supplies a radical extension. Its group is therefore solvable.
\item False. Although both are transcendental, they satisfy the polynomial relation $Y^2-X=0$ over $\Q$, so they are algebraically dependent.
\item True. The group is cyclic, hence abelian and solvable.
\item False. $x^5-2$ is solved by adjoining a fifth root of 2 and a fifth root of unity.
\item True. The quartic formulas, or solvability of subgroups of $S_4$, prove it.
\item False. A real cube root of 2 has degree 3 over $\Q$ and cannot belong to a tower of quadratic extensions, whose degree is a power of 2.
\item False. The cubic $x^3-x$ has rational roots $-1,0,1$, requiring no extension at all.
\item True. A subnormal series with abelian quotients characterizes solvable groups, linking the group criterion to radical towers.
\end{enumerate}

\exercise{56.4}{The quadratic formula over a field}
Multiplying $ax^2+bx+c=0$ by $4a$ and completing the square gives $(2ax+b)^2=b^2-4ac$. Since $a\ne0$ and characteristic is not 2, $2a$ is invertible. With $\delta^2=b^2-4ac$, the roots are $(-b+\delta)/(2a)$ and $(-b-\delta)/(2a)$, so they lie in $F(\delta)$. Conversely either root determines $\delta=2a\alpha+b$, up to sign, so the field generated by the roots contains $\delta$. Thus the splitting field is $F(\sqrt{b^2-4ac})$, including the trivial-extension cases where the discriminant is already a square or zero.

\exercise{56.5}{A biquadratic solved by successive radicals}
Put $y=x^2$. First solve $ay^2+by+c=0$ by adjoining $\delta=\sqrt{b^2-4ac}$, giving $y_\pm=(-b\pm\delta)/(2a)$. Next adjoin $\sqrt{y_+}$ and $\sqrt{y_-}$. The resulting tower is radical, and contains all four roots $\pm\sqrt{y_+},\pm\sqrt{y_-}$, with coincidences allowed. Hence the splitting field is contained in this radical extension. The characteristic assumption permits division by $2a$.

\exercise{56.6}{Refining a series with abelian factors}
Consider one original step $N\trianglelefteq M$ with $M/N$ abelian, and a refined chain $N=H_0\trianglelefteq H_1\trianglelefteq\cdots\trianglelefteq H_k=M$. For each $j$, $H_j/N$ is a subgroup of the abelian group $M/N$, hence is abelian. By the third isomorphism theorem,
\[H_j/H_{j-1}\cong(H_j/N)/(H_{j-1}/N),\]
which is a quotient of an abelian group and is therefore abelian. Every new step lies within some original step, so this proves the assertion for the whole refinement.

\exercise{56.7}{Refining solvable factors to abelian composition factors}
For a step $N\trianglelefteq M$ whose quotient $M/N$ is solvable, choose a subnormal series of $M/N$ with abelian factors. Lift its subgroups under the quotient map $M\to M/N$; the correspondence and isomorphism theorems give a subnormal refinement from $N$ to $M$ with those same abelian quotients. Do this for every original step. Because the group is finite, refine further until every nontrivial quotient is simple, obtaining a composition series. Exercise 56.6 ensures the refined quotients remain abelian. Thus all composition factors are simple abelian finite groups, hence cyclic groups of prime order.

\exercise{56.8}{An explicit quintic with no radical solution}
\begin{enumerate}[label=(\alph*)]
\item Label the letters by $\Z_5$ so the five-cycle is $c:j\mapsto j+1$. If the transposition is $(u,v)$, its conjugates $c^k(u,v)c^{-k}=(u+k,v+k)$ all lie in $H$. Regard these as edges joining vertices differing by the nonzero residue $v-u$. Such edges make a connected graph, since every nonzero residue generates the additive group $\Z_5$. Edge transpositions of a connected graph generate all transpositions: along a path, conjugating adjacent edge swaps gives a swap of the path endpoints, by induction on its length. Hence $H$ contains every transposition and therefore $H=S_5$.
\item Irreducibility of the quintic makes the Galois action on its five roots transitive, so the orbit-stabilizer theorem gives $5\mid|G|$. Cauchy's theorem gives an element of order 5, which in $S_5$ is a five-cycle. Complex conjugation exchanges exactly the two nonreal roots and fixes the three real ones, so its permutation is a transposition, not merely an arbitrary element of order 2. Part (a) now gives $G=S_5$.
\item For $f(x)=2x^5-5x^4+5$, Eisenstein at 5 proves irreducibility: 5 divides every nonleading coefficient, 25 does not divide the constant term, and 5 does not divide 2. The derivative is $f'(x)=10x^3(x-2)$. Thus $f$ increases on $(-\infty,0)$, decreases on $(0,2)$, and increases on $(2,\infty)$. Its local maximum is $f(0)=5>0$ and its local minimum is $f(2)=-11<0$. The odd degree with positive leading coefficient gives limits $-\infty$ and $+\infty$ at the two ends. Therefore continuity and strict monotonicity give exactly one real zero on each of these three intervals, and no others. Its remaining two roots are a nonreal conjugate pair. Here is a sketch, with the two extrema labeled and the curve drawn from actual sampled values; arrows indicate the tails.
\begin{center}\setlength{\unitlength}{1pt}
\begin{picture}(360,150)
\put(0,70){\vector(1,0){350}}\put(108,0){\vector(0,1){143}}
\put(352,65){\makebox(0,0){$x$}}\put(100,145){\makebox(0,0){$f(x)$}}
\put(108,90){\circle*{3}}\put(288,26){\circle*{3}}
\put(144,105){\makebox(0,0){$(0,5)$}}\put(275,11){\makebox(0,0){$(2,-11)$}}
\polyline(0.00,28.62)(9.00,47.83)(18.00,62.00)(27.00,72.15)(36.00,79.19)(45.00,83.85)(54.00,86.79)(63.00,88.50)(72.00,89.41)(81.00,89.82)(90.00,89.97)(99.00,90.00)(108.00,90.00)(117.00,90.00)(126.00,89.97)(135.00,89.86)(144.00,89.57)(153.00,89.00)(162.00,88.03)(171.00,86.54)(180.00,84.43)(189.00,81.60)(198.00,78.00)(207.00,73.60)(216.00,68.43)(225.00,62.58)(234.00,56.19)(243.00,49.50)(252.00,42.81)(261.00,36.55)(270.00,31.21)(279.00,27.45)(288.00,26.00)(297.00,27.77)(306.00,33.78)(315.00,45.23)(324.00,63.46)(333.00,90.00)(342.00,126.56)
\put(2,26){\vector(-1,-2){4}}\put(340,126){\vector(1,2){5}}
\end{picture}\end{center}
By (b) its Galois group is $S_5$. This group is not solvable: its simple nonabelian normal subgroup $A_5$ obstructs an abelian-factor series. The solvability theorem therefore shows this quintic is not solvable by radicals over $\Q$.
\end{enumerate}
